\section{Limiti del lavoro}

OMNeT++, SUMO, PLEXE e Simu5G permettono uno studio utile e coerente, ma non descrivono la complessità del mondo reale in modo completo. Di conseguenza, i risultati vanno letti soprattutto come validazione funzionale del sistema, non come previsione diretta delle prestazioni su strada. Questo vale in particolare per il confronto tra LTE e 5G, che dipende anche dai modelli radio e dai parametri di handover adottati.

Una seconda limitazione è l'omogeneità dello scenario. I veicoli del platoon sono tutti uguali per modello dinamico, controllore e dotazione di comunicazione, mentre nella realtà è ragionevole aspettarsi differenze tra mezzi, sensori, antenne e tempi di attuazione. Anche questo può influenzare sia la stabilità del platoon sia l'efficacia del \emph{fallback}.

Infine, il canale di comunicazione è stato modellato in modo selettivo. Lo scenario realistico include interferenze e handover, ma non considera in modo completo tutte le possibili fonti di degrado, come fading e shadowing, perdità di linea di vista, traffico radio esterno o condizioni ambientali variabili. Per questo la tesi non fornisce uno scenario definitivo, ma una base coerente per analisi successive più realistiche.
